\scanpage{15}
En effet, \unsure{évidemt|évident|évidemment} $X'\ne\varnothing$. Soit
$X'_0$ une composante de dim \unsure{minima|minimum},
soit $l$, \inserted{donc $l\geq0$} prouvons $l\unsure{\geq|>} k+1$. \deleted{Sinon il
existera $f_1\cdots f_l$} \deleted{Si on avait $l\leq k$,
alors $X'$ étant connexe [prendre des
$f_i=0$] on \unsure{conclurait|concluerait} que $X'=X'_0$, donc
$X'$ de dim 0, mais alors $\exists\,f_1\in rA$
tel que $Y'=X'\cap V(f)=\varnothing$, \deleted{\unsure{à fortiori}}\inserted{\unsure{que}}
$X'\cap\underbrace{V(f)\cap\cdots\cap V(f)}_{k\ \text{fois}}=\varnothing$, absurde.}
\deleted{Si on avait $l\leq k$,} il existe\deleted{\unsure{ra}}
$l$ \unsure{élémts|éléments} de $rA$ \deleted{tels que} $f_1\cdots f_l$
tels que, \inserted{\unsure{fixant} \deleted{\unsure{pour|tous}} \unsure{les|tte|toute} composantes irréductibles $X'_i$, de dim $n_i$,}
$V(f_1)\cap\cdots\cap \unsure{V|N}(f_l)\cap X'_{\deleted{\unsure{o}}i}\unsure{(=}Y'_i)$ \deleted{\unsure{soit de dim}}
\inserted{les \unsure{comp. irréd.} de $Y'_i$} \deleted{ait \unsure{tt. comp.}} de dim $n_i-l$, et \deleted{\unsure{alors ces composantes}} irréductibles
\deleted{\unsure{fermés} \unsure{et} \unsure{que} \unsure{les} \unsure{pts} \unsure{soient} \unsure{clairs}} sont \unsure{de} distinctes ($i$ variable).

\deleted{\unsure{de ttes} composantes irréd. de $X'\cap V(f_1)\cap\cdots\cap V(f_l)$.
Comme ce dernier est connexe, il
serait de dim 0 et \unsure{réduit à}
\unsure{un} pt \unsure{Si on avait} $l\leq k$,}

En particulier, pour $i=0$, $Y'_i$ est fini,
donc les pts de $Y'_i$ sont des composantes
irréductibles de dim 0. \inserted{\uuline{Si} $l\leq k$,} l'hyp.
implique que $Y'$ est connexe, donc
$X'$ serait irréductible de dim $l\leq k$
\deleted{Si $l<k$, \unsure{absurde}, \unsure{car} une \unsure{équation|section} de}
% archive sheet number: 14
